\subsection{Two spin endpoints}\label{p10-sec:examples}

Let $A,B$ be distinct theta characteristics on $C$ with
$H^0(C,A)=H^0(C,B)=0$. For a theta characteristic $S$, take the
positive spin operator $D_S$ associated with a lift of the
uniformizing representation to $\mathrm{SL}(2,\RR)$. In a spin frame
on the upper half-plane it is given by
\[
                D_S u=\bar\partial u+
                    \frac{i\bar u}{2\operatorname{Im}z}\,d\bar z .
\]
The choice of lift specifies the signs in the transition functions of
the frame. On the same underlying smooth complex rank-two bundle
consider the ordered endpoints
\[
               D_-=D_A\oplus D_A,\qquad D_+=D_A\oplus D_B .
\]
They are rigid in every degree. Indeed, on a finite cover $f$, the
equation for one spin summand makes $d(u^2 f^*dz)$ a nonnegative area
form whose density is positive wherever $u\ne0$ and $df\ne0$.
Integration forces $u=0$ on every positive-degree component, and nodal
matching gives vanishing on each attached contracted subcurve. This
also proves base injectivity when the antilinear term is multiplied by
any positive scalar. At scalar zero each base Dolbeault operator is
invertible by non-effectivity and Riemann--Roch. Hence the two endpoint
determinants have sign $\epsilon=+1$. To connect the operators in the
base-invertible locus, turn off the antilinear term on the second
summand, join the two Dolbeault operators in the complement of the
theta divisor in $\operatorname{Pic}^{g-1}(C)$, and restore the term
on $B$. The complement is path connected since the divisor has real
codimension two. Trivialize the smooth bundles along the interval and
perturb the path relative to its endpoints.

We evaluate the two endpoints using the absolute partition formula
of \cite[Theorem~1.1]{WueSpin}. Write $c=g-1$ and
$X_k=(2\sin(kv/2))^2$. At the repeated spin the disconnected
degree-$d$ coefficient is the sum over partitions $\lambda\vdash d$
with hook weight $\prod_{a\in\lambda}X_{h(a)}^c$. At the distinct
pair with positive determinant the same weights are summed over just
the self-conjugate partitions. These two absolute evaluations are
the input from spin theory. The signed number of crossings between
them follows from Theorem~\ref{p10-thm:chambers}.

Let $Z_d^-$ and $Z_d^+$ be the respective disconnected coefficients,
requiring positive degree on every connected component. For $d\le3$
the partition sums read
\[
\begin{array}{c|cc}
 d & Z_d^- & Z_d^+ \\ \hline
 1 & x^c & x^c \\
 2 & 2(xX_2)^c & 0 \\
 3 & 2(xX_2X_3)^c+(x^2X_3)^c & (x^2X_3)^c .
\end{array}
\]
Indeed the row and column of size three have hooks $1,2,3$, while
$(2,1)$ has hooks $1,1,3$. The logarithm of the disconnected series
has degree-two coefficient $Z_2-Z_1^2/2$ and degree-three coefficient
$Z_3-Z_1Z_2+Z_1^3/3$. Using $X_2=x(4-x)$ and $X_3=x(3-x)^2$ gives
the connected differences in the order $D_+$ minus $D_-$,
\begin{align*}
       x^{-2c}\Delta L_2 &=-2(4-x)^c,\\
       x^{-3c}\Delta L_3 &=2(4-x)^c\bigl(1-(3-x)^{2c}\bigr).
\end{align*}

\begin{corollary}\label{p10-cor:spin-counts}
For every generic path between these ordered endpoints in the
base-invertible locus, the signed polynomial of smooth cubic crossings is
\begin{equation}\label{p10-eq:spin-crossings}
       W_3(x)=\bigl(4-x(1-x)^2\bigr)^c
                         -\bigl((4-x)(3-x)^2\bigr)^c .
\end{equation}
The largest normalization genus with nonzero signed count is $6c$.
Every such path contains at least $8c$ smooth cubic cover--parameter
pairs in that genus.
\end{corollary}

\begin{proof}
Insert $W_2(x)=-2(4-x)^c$ and the cubic difference into
\eqref{p10-eq:chambers}. The terms in $(4-x)^c$ cancel, giving
\eqref{p10-eq:spin-crossings}. The two cubics on the right, before taking
$c$th powers, have common leading coefficient $-1$ and difference
$-8(x-2)^2$. Their difference of $c$th powers thus has degree $3c-1$
and leading coefficient $-8c(-1)^{c-1}$. By Riemann--Hurwitz, a
cubic cover with $2l$ simple branch values has genus $1+3c+l$.
Taking $l=3c-1$ proves both assertions.
\end{proof}

On a genus-two target one can take any two distinct even theta
characteristics. Each has degree one, hence at most one independent
section, and even parity forces non-effectivity. Here
$W_2(x)=-8+2x$ and $W_3(x)=-32+32x-8x^2$. The signed counts of smooth
cubics are $-32$, $+32$, and $-8$ in normalization genera four, five,
and six, respectively. In particular the half-coefficient of $x^2$ in
$x^{-3}\Delta L_3$ is $-10$, and subtracting the double-cover term in
\eqref{p10-eq:chambers} raises the signed four-branch count to $-8$.
Every such path therefore contains at least $72$ smooth cubic
cover--parameter pairs across these three genera. In higher genera
the signed count vanishes.
